\bookchapter{Capitolo 41: Panorama degli strumenti di programmazione AI nel 2026}

\chapterintro{Una prevenzione prima di iniziare: questo capitolo è uno “scatto” dell'ecosistema degli strumenti di agosto 2026. Tra sei mesi sarà superato — ed è proprio la prima lezione che il capitolo vuole insegnarti. Lo scopo del capitolo non è insegnarti “quale strumento usare”, ma aiutarti a costruire un quadro di giudizio su “come scegliere gli strumenti”. Quando capirai la filosofia di design dei diversi strumenti, potrai in ogni momento fare la scelta giusta per te.}

\booksection{Le tre scuole e il quadro di scelta}

Nel 2026 gli strumenti di programmazione AI si dividono all'incirca in tre grandi scuole, ognuna basata su una filosofia di design diversa. Prima di tutto, guarda questo panorama:

\begin{center}
\begin{tabularx}{\linewidth}{|X|X|X|X|}
\hline
\textbf{Scuola} & \textbf{Strumenti rappresentativi} & \textbf{Filosofia di design} & \textbf{Scenari adatti} \tabularnewline
\hline
AI integrata nell'IDE & Cursor 3, Trae, Windsurf, CodeBuddy & Integrare in profondità le capacità AI nell'editor & Sviluppo quotidiano, code review, potenziamento graduale \tabularnewline
Agent da terminale & Claude Code, Codex CLI, Aider & Il terminale come interfaccia: l'AI manipola direttamente file system e comandi & Operazioni a livello di progetto, modifiche in blocco, integrazione CI/CD \tabularnewline
Builder nel cloud & v0, Bolt, Replit Agent & Creare da zero un progetto completo nel browser & Prototipi rapidi, partenza senza configurazione, anteprima immediata \tabularnewline
\hline
\end{tabularx}
\end{center}

\begin{center}\small\color{gray}Le tre scuole degli strumenti di programmazione AI\end{center}

La chiave per scegliere uno strumento non è “quale è il migliore”, ma “quale è più adatto al tuo scenario attuale”. Tre dimensioni di giudizio pratiche:

\begin{itemize}
  \item Qual è la tua base di programmazione? Da zero scegli un Builder nel cloud, con una base scegli l'AI integrata nell'IDE, per utenti avanzati l'Agent da terminale.
  \item A che fase è il tuo progetto? Partendo da zero scegli il Builder nel cloud, per lo sviluppo iterativo l'AI integrata nell'IDE, per manutenzione e refactoring l'Agent da terminale.
  \item Quali sono le tue esigenze di deployment? Per una verifica rapida il Builder nel cloud, per i progetti seri l'AI integrata nell'IDE, per i flussi automatizzati l'Agent da terminale.
\end{itemize}

\booksection{La scuola dell'AI integrata nell'IDE}

L'AI integrata nell'IDE è la scuola più diffusa: incorpora le capacità AI nel tuo editor di codice, mantenendo l'esperienza tradizionale di programmazione e aggiungendo il potenziamento AI. Nel 2026 i rappresentanti di questa scuola sono Cursor 3, Trae, Windsurf e CodeBuddy.

\booksubsection{Cursor 3}

Cursor è attualmente l'IDE di programmazione AI più usato al mondo. Cursor 3, rilasciato nell'aprile 2026, segna una svolta importante per gli strumenti di programmazione AI: l'utilizzo della modalità Agent ha superato per la prima volta quello del completamento Tab — gli sviluppatori sono passati da “l'AI mi completa il codice” a “l'AI completa i miei compiti”. Le caratteristiche principali di Cursor 3 includono: editing multi-file (l'AI può modificare più file insieme per portare a termine il compito), consapevolezza del contesto (comprende automaticamente struttura del progetto e dipendenze) e profonda integrazione con il terminale (l'AI può eseguire comandi e vedere l'output direttamente dall'editor).

\booksubsection{Trae}

Trae è l'IDE di programmazione AI lanciato da ByteDance, con una base di utenti molto ampia nel mercato cinese. La sua modalità SOLO (chiamata anche modalità Builder) è l'archetipo della capacità “creare un progetto con una frase”: tu descrivi il requisito e l'AI completa in autonomia tutto il flusso, dall'inizializzazione del progetto alla generazione del codice. I punti di forza di Trae sono il supporto al cinese e l'adattamento alla rete nazionale, perfetti per gli utenti cinesi che vogliono iniziare in fretta.

\booksubsection{Altri strumenti}

Windsurf (ex Codeium) è noto per un'esperienza di completamento del codice fluidissima, ideale per chi cerca la sensazione di “scrivere al volo”. CodeBuddy di Tencent Cloud, come estensione di VS Code, eccelle nell'integrazione con l'ecosistema cloud ed è adatto a progetti pensati per Tencent Cloud o per lo sviluppo di Mini Program WeChat. Anche Lingma di Alibaba (Tongyi Lingma) e Comate di Baidu offrono esperienze di integrazione nell'IDE simili, ciascuna con ottimizzazioni proprie per gli scenari in cinese.

\booksection{La scuola degli Agent da terminale}

L'Agent da terminale è la scuola che cresce più in fretta nel 2026. Idea centrale: l'AI non dovrebbe essere un semplice plugin dell'editor, ma poter operare sull'intero ambiente di sviluppo come fa uno sviluppatore — leggere e scrivere file, eseguire comandi, vedere l'output, installare dipendenze. Claude Code è il rappresentante di questa scuola.

\booksubsection{Claude Code}

Claude Code è lo strumento di programmazione AI da terminale lanciato da Anthropic, arrivato a metà 2026 a un fatturato annualizzato di 2,5 miliardi di dollari. Il suo funzionamento è completamente diverso da quello di un IDE tradizionale: descrivi il compito in linguaggio naturale nel terminale, e Claude Code legge in autonomia i file pertinenti, modifica il codice, esegue i test, controlla i risultati e corregge iterando — senza che tu debba aprire l'editor.

Il vantaggio di Claude Code è la “percezione dell'intero progetto”: non è limitato dalla finestra dell'editor, può esplorare liberamente l'intera directory del progetto e capire le dipendenze tra i file. Supporta anche i file di memoria del progetto CLAUDE.md, così l'AI conosce il contesto fin dall'inizio di ogni sessione. Per refactoring di progetti grandi, modifiche in blocco, scrittura di test e compiti simili, un Agent da terminale è più efficiente di un plugin per IDE.

\booksubsection{Codex CLI e Aider}

Codex CLI di OpenAI offre un'esperienza di programmazione da terminale simile e supporta la funzione /import — permette di importare una base di codice esistente perché l'AI la comprenda. Aider è uno strumento open source di programmazione AI da terminale che supporta più modelli (GPT, Claude, DeepSeek e altri), adatto a chi vuole cambiare modello con flessibilità. Questi strumenti hanno ognuno le sue peculiarità, ma la stessa idea di fondo: dare all'AI, nel terminale, la piena capacità di operare sull'ambiente di sviluppo.

\booksection{La scuola dei Builder nel cloud}

Il Builder nel cloud è la scuola più adatta a chi parte da zero. Non devi installare nulla: apri il browser, descrivi il tuo requisito, e l'AI crea nel cloud un progetto completo con anteprima immediata.

\begin{itemize}
  \item v0 (di Vercel): bravo a generare componenti e pagine UI per il frontend, ideale per realizzare prototipi web in fretta.
  \item Bolt: crea applicazioni full-stack partendo da una sola frase, con frontend, backend e database.
  \item Replit Agent: lavora nell'IDE cloud di Replit, non solo crea progetti ma li esegue e li distribuisce direttamente nel cloud.
\end{itemize}

Il vantaggio comune di questi strumenti è “iniziare senza configurazione” — apri il browser e via. Il limite: quando il progetto si complica, la flessibilità dell'ambiente di editing nel cloud non eguaglia quella di un IDE locale. Per questo sono più adatti alla fase prototipale “da 0 a 1”: verifica rapidamente l'idea con un Builder nel cloud e, una volta validata, migra su un IDE locale per lo sviluppo approfondito.

\booksection{Gli strumenti cambiano, la metodologia no}

Se domani tutti gli strumenti di programmazione con l'IA sparissero, dovrebbero restare tre abitudini: pensare bene prima di agire, fare un piccolo passo alla volta, verificare una volta finito. Gli strumenti sono amplificatori, la metodologia è la fondazione. Gli amplificatori si possono sostituire, la fondazione non deve vacillare.

Ripercorriamo i cambiamenti di strumenti dal 2024 al 2026: nel 2024 il re era GitHub Copilot, nel 2025 lo ha sostituito Cursor, nel 2026 Claude Code ha fatto il grande balzo. A ogni nuovo strumento qualcuno proclama “stavolta è diverso”, ma alla fine si scopre che chi sfrutta bene lo strumento nuovo è proprio chi ha una metodologia solida. A loro basta qualche giorno di adattamento per cambiare strumento, mentre chi ha imparato solo i comandi riparte da zero a ogni cambio.

Perciò il percorso di apprendimento consigliato da questo libro è: inizia con uno strumento (uno qualsiasi va bene) e assorbi nella pratica i cinque principi di pensiero; quando la tua metodologia sarà abbastanza solida, prova strumenti di scuole diverse e senti le differenze di filosofia di design; alla fine formerai la tua combinazione personale — per esempio Cursor per lo sviluppo quotidiano, Claude Code per refactoring e test, v0 per i prototipi rapidi. La combinazione varia da persona a persona, ma la metodologia è comune.

\begin{tcbpractice}{Mettiti alla prova: prepara la tua “combinazione di armi”}
Posizionati con le tre dimensioni di giudizio (base di programmazione, fase del progetto, esigenze di deployment): se parti da zero prova Bolt o v0 per fare una pagina; se hai una base fai un compito iterativo con Cursor o Trae; se vuoi una sfida esegui Claude Code una volta nel terminale. Poi rispondi: quale scuola rispecchia meglio la tua fase attuale?

\end{tcbpractice}

\begin{tcbsummary}{Riepilogo del capitolo}
\begin{itemize}
  \item Le tre scuole: AI integrata nell'IDE (sviluppo quotidiano), Agent da terminale (operazioni a livello di progetto), Builder nel cloud (prototipi da zero)
  \item Per scegliere lo strumento, tre domande: base di programmazione? fase del progetto? esigenze di deployment?
  \item In Cursor 3 la modalità Agent ha superato per la prima volta il completamento Tab; Claude Code fattura 2,5 miliardi di dollari annualizzati
  \item I Builder nel cloud sono ideali per il “da 0 a 1”: dopo la verifica, si torna a un IDE locale per lo sviluppo approfondito
  \item Gli strumenti si rinnovano ogni tre mesi: chi ha imparato la metodologia cambia strumento con pochi giorni di adattamento
\end{itemize}

\end{tcbsummary}
